\anexo{Preguntas de la Encuesta}

\section{Preguntas de la Encuesta}

\subsection{Sección 1: Filtro y Perfilado}

Esta sección fue dirigida a todos los encuestados. Al finalizarla, el usuario avanza según lo que responda en la pregunta Q2.

\begin{enumerate}[label=\textbf{Q\arabic*:}, leftmargin=*]
\item ¿Buscaste una propiedad en CABA en los últimos 10 años o estás buscando una actualmente?

\begin{itemize}
\item Sí: continúa.
\item No: enviar formulario.
\end{itemize}

\item ¿En qué rango de edad te encontrás?

\begin{itemize}
\item 18--25.
\item 26--35.
\item 36--45.
\item 46--55.
\item 56+.
\end{itemize}

\item ¿En qué barrio/s de CABA buscaste principalmente?

Tipo de respuesta: menú desplegable. Se incluyó la lista oficial de los 48 barrios de CABA ordenados alfabéticamente, desde Agronomía, Almagro y Balvanera hasta Villa Urquiza.

\item ¿Qué tipo de propiedad estabas buscando principalmente?

\begin{itemize}
\item Departamento.
\item Casa.
\item PH (Propiedad Horizontal).
\item Múltiples opciones / Indistinto.
\end{itemize}

\item ¿Qué portales inmobiliarios utilizaste habitualmente?

Casillas de verificación con posibilidad de elegir varias opciones.

\begin{itemize}
\item Zonaprop.
\item Argenprop.
\item Mercado Libre Inmuebles.
\item Cabaprop / Mudafy / Remax (portales de franquicias).
\item Grupos de Facebook o redes sociales.
\item Inmobiliarias tradicionales (sitio web propio o presencial).
\item Otro.
\end{itemize}

\item ¿Qué tipo de operación estabas buscando?

\begin{itemize}
\item Opción A: Alquiler $\rightarrow$ ir a Sección 2A.
\item Opción B: Compra $\rightarrow$ ir a Sección 2B.
\end{itemize}
\end{enumerate}

\subsection{Sección 2A: El Dolor del Precio}

Rama exclusiva para alquiler. Aquí se ajustó la redacción para el inquilino. Al finalizar, el usuario avanza a la Sección 3.

\begin{enumerate}[label=\textbf{Q\arabic*:}, leftmargin=*, start=7]
\item Del 1 al 5, ¿qué tan difícil te resulta saber si el precio mensual del alquiler que piden es justo o está fuera de mercado?

Escala lineal: 1 (Muy fácil) a 5 (Muy difícil).

\item Cuando ves un alquiler publicado, ¿qué tanta desconfianza te genera el valor inicial en pesos respecto al estado real del inmueble?

Escala lineal: 1 (Nada de desconfianza, confío en el precio) a 5 (Mucha desconfianza, asumo que está inflado).

\item Actualmente, ¿cómo intentás averiguar si el precio del alquiler es adecuado?

Opción múltiple.

\begin{itemize}
\item Comparo ``a ojo'' con otras publicaciones similares en portales.
\item Le consulto a familiares, amigos o conocidos.
\item Me guío por lo que me dice el corredor inmobiliario.
\item No tengo un método claro, me guío únicamente por mi presupuesto disponible.
\item Otro.
\end{itemize}

\item ¿Alguna vez sentiste que terminaste pagando (o te pidieron) un alquiler más caro de lo que realmente valía el departamento?

\begin{itemize}
\item Sí, definitivamente.
\item Sí, pero lo acepté por la urgencia/necesidad de mudarme.
\item No, creo que siempre me pidieron precios justos.
\item No estoy seguro/a.
\end{itemize}
\end{enumerate}

\subsection{Sección 2B: El Dolor del Precio}

Rama exclusiva para compra. Aquí se ajustó la redacción para el comprador. Al finalizar, el usuario avanza a la Sección 3.

\begin{enumerate}[label=\textbf{Q\arabic*:}, leftmargin=*, start=7]
\item Del 1 al 5, ¿qué tan difícil te resulta saber si el valor de venta (en dólares) de una propiedad es justo o está sobrevaluado?

Escala lineal: 1 (Muy fácil) a 5 (Muy difícil).

\item Cuando ves una propiedad en venta, ¿qué tanta desconfianza te genera el valor publicado en dólares respecto al estado real del inmueble?

Escala lineal: 1 (Nada de desconfianza, confío en el precio) a 5 (Mucha desconfianza, asumo que está inflado).

\item Actualmente, ¿cómo intentás averiguar si el precio de publicación es adecuado para hacer una oferta?

Opción múltiple.

\begin{itemize}
\item Comparo manualmente con otras propiedades similares en la misma zona.
\item Pido tasaciones a otros profesionales o consulto con un corredor de confianza.
\item Busco datos históricos del valor del m$^2$ en el barrio.
\item Voy a visitarla y oferto agresivamente hacia abajo ``por las dudas''.
\item Otro.
\end{itemize}

\item ¿Alguna vez descartaste ir a visitar una propiedad porque sentiste que el precio en dólares era un ``disparate'' sin justificación?

\begin{itemize}
\item Sí, me pasa frecuentemente.
\item Sí, me pasó alguna vez.
\item No, si me gusta voy a verla e intento negociar el precio cara a cara.
\end{itemize}
\end{enumerate}

\subsection{Sección 3: El Dolor del Entorno}

En esta sección las ramas se unen nuevamente. Está dirigida tanto a compradores como a inquilinos, ya que ambos perfiles enfrentan problemas similares vinculados con el entorno urbano y la falta de información.

\begin{enumerate}[label=\textbf{Q\arabic*:}, leftmargin=*, start=11]
\item Cuando encontrás una propiedad que te gusta, ¿qué acciones extra realizás para investigar la zona?

\begin{itemize}
\item Abro Google Maps para ver qué comercios o transportes hay cerca.
\item Uso Google Street View para ver la fachada, la cuadra y el estado de las veredas.
\item Voy a caminar por el barrio en distintos horarios (por ejemplo, de noche).
\item Busco noticias u opiniones sobre la seguridad en esa zona específica.
\item No hago nada de esto, confío en la descripción del aviso.
\end{itemize}

\item Del 1 al 5, ¿qué nivel de importancia le das a tener cerca los siguientes puntos de interés (POIs)?

\begin{center}
\scriptsize
\begin{longtable}{|p{0.30\textwidth}|*{5}{>{\centering\arraybackslash}p{0.095\textwidth}|}}
\caption{Cuadrícula de importancia de puntos de interés}\label{tab:pois-encuesta}\\
\hline
 & 1 & 2 & 3 & 4 & 5 \\
\hline
\endfirsthead
\hline
 & 1 & 2 & 3 & 4 & 5 \\
\hline
\endhead
Transporte público (Subte / Metrobus / Tren) &  &  &  &  &  \\
\hline
Supermercados y comercios esenciales &  &  &  &  &  \\
\hline
Parques, plazas o espacios verdes &  &  &  &  &  \\
\hline
Centros de salud (Hospitales / Clínicas) &  &  &  &  &  \\
\hline
Polos gastronómicos o vida nocturna &  &  &  &  &  \\
\hline
Instituciones educativas (Facultades / Colegios) &  &  &  &  &  \\
\hline
\end{longtable}
\end{center}

\item El ruido ambiental (tráfico, avenidas, bares), ¿alguna vez te hizo descartar una propiedad o arrepentirte luego de mudarte?

\begin{itemize}
\item Sí, descarté propiedades antes de mudarme al notar que era una zona ruidosa.
\item Sí, me mudé y me arrepentí por ruidos constantes que no había notado al visitarla.
\item Me molesta el ruido, pero no al punto de descartar una propiedad o arrepentirme.
\item No me importa el ruido ambiental.
\end{itemize}

\item ¿Cuál es tu mayor ``Dealbreaker'' (motivo de descarte inmediato) respecto al entorno de una propiedad?

\begin{itemize}
\item Falta de conexión rápida con transporte público.
\item Zona ruidosa (avenidas muy transitadas, tren, bares nocturnos).
\item Sensación de inseguridad en las calles aledañas.
\item Falta de comercios esenciales para el día a día.
\item Lejanía de plazas o espacios verdes.
\item Otro.
\end{itemize}
\end{enumerate}

\subsection{Sección 4: Validación de la Solución Nidio}

Esta sección fue dirigida a todos los encuestados.

\begin{enumerate}[label=\textbf{Q\arabic*:}, leftmargin=*, start=15]
\item Si un portal te indicara objetivamente un ``rango de precio esperado'' basado en datos reales, ¿qué tanto influiría en tu decisión?

Escala lineal: 1 (No influiría nada, confío solo en mi criterio) a 5 (Sería totalmente decisivo).

\item Si la plataforma te ofreciera un ``Life Score'' personalizado sobre qué tan compatible es la propiedad elegida con tus intereses (cercanía a transporte, salud, espacios verdes, comercios, etc.), ¿lo considerarías útil?

\begin{itemize}
\item Muy útil, me ahorraría el trabajo de buscar manualmente en los mapas.
\item Algo útil, lo miraría como referencia pero igual haría mi propia investigación.
\item Poco útil, prefiero analizar el mapa de la zona 100\% por mi cuenta.
\end{itemize}

\item Para finalizar, si una nueva plataforma te ofreciera todas estas herramientas de análisis integradas, ¿cambiarías tu forma actual de buscar propiedades?

\begin{itemize}
\item Sí, la elegiría como mi portal principal en lugar de los tradicionales.
\item La usaría a la par de los portales de siempre (Zonaprop, Argenprop, etc.) para comparar.
\item No, prefiero seguir usando únicamente los portales tradicionales que ya conozco.
\end{itemize}
\end{enumerate}
